% 本文件由 LaTeX 排版 Agent 根据有效 translation.json 自动生成。
% author=John Chrysostom; work_id=1912; title=若你的仇人饿了，你要给他吃

\chapter{\Quote{若你的仇人饿了，你要给他吃}}

% structure: p0001 source=h1 target=omitted reason=page-title-already-rendered-by-parent

\Emph{致那些没有参加集会的人。}\par

致那些没有参加集会的人；论宗徒的这句话：\Quote{若你的仇人饿了，你要给他吃，}以及对伤害的怨恨。\par

1. 看来我最近向你们发表的那篇冗长演说，旨在激发你们对这里集会的热忱，并没有产生什么好处：因为我们的教会再次缺少她的儿女。为此，我又被迫显得惹人厌烦和负担重重，责备那些在场的人，责怪那些缺席的人：对他们，是因为他们没有除去懒惰；对你们，是因为你们没有为弟兄们的救恩伸出援手。我被迫显得负担重重、惹人厌烦，不是为了我自己，也不是为了我自己的财产，而是为了你们和你们的救恩——这对我来说比任何东西都更宝贵。谁若愿意，尽管生气，称我傲慢无礼，我仍不会停止为同一目的不断烦扰他；因为对我来说，再没有比这种厚颜更有益的了。因为或许，或许，至少这一点能使你们感到羞愧，为了避免因同一件事被反复纠缠，你们会参与到对弟兄们的关怀中。当我没看到你们在德行上进步时，称赞对我有何益处？当我见到你们的敬虔增长时，听众的沉默又有何害处？因为说话者的称赞不在于掌声，而在于听众对虔诚的热忱；不在于聆听时发出的喧哗，而在于持久的认真。掌声一从口中发出，就在空中消散、灭亡；但听众品德的改善，给说话者和听从者都带来不朽和永生的赏报。对你们喝彩的称赞使说话者在世上显赫，但你们灵魂的敬虔在基督的审判台前给教师带来极大的信心。因此，如果有人爱那说话者，就不要渴望听众的喝彩，而要渴望他们的益处。忽视我们的弟兄并不是普通的过错，而是招致严厉惩罚和不可饶恕的刑罚的。埋藏塔冷通的比喻证明了这一点：那个人至少没有因自己的品行被责备，因为就那笔存款本身而言，他并不是一个坏人，因为他原封不动归还了；但他在管理存款方面却是个坏人，因为他没有使托付给他的钱倍增，因而受了惩罚。由此可见，即使我们热心、训练有素、对聆听圣经大有热情，这也不足以使我们获得救恩。因为那笔存款必须倍增，当我们除了自己的救恩之外，也为他人的益处作些安排时，它就倍增了。因为比喻中的那个人说：\Quote{看，你所有的都在这里。}但这并不能为他辩护，因为对他说：\Quote{你应当把我的银子交给兑换银钱的人。} \BibleRef{玛 25:27}\par

且请注意，主人的命令是何等容易：因为人们确实让那些放贷收利息的人负责追回款项；有人说：\Quote{你存了款，你必须收回它；我与收款人无关。}但天主不是这样做的；他只命令我们存放，并不要求我们追回。因为说话者有权劝告，却无权说服。因此他说：\Quote{我只让你负责存放，而不负责追回。}还有什么比这更容易的呢？可是那仆人竟称如此温柔仁慈的主人为苛刻的。因为忘恩负义和懒惰的人常有这种习惯：他们总是试图把自己的过错归咎于主人。因此，那人被连拖带绑扔到外面的黑暗里。为了使我们免于这种惩罚，让我们把教导交给弟兄们，无论他们是否被说服。因为他们若被说服，就会使自己和我们都得益处；若不被说服，他们固然使自己陷入不可避免的惩罚，但绝不能对我们造成丝毫伤害。因为我们已尽了自己的本分，劝告了他们；若他们不听，我们不会因此受害。因为责备将落在我们身上，不是因为我们未能说服，而是因为我们未能劝告；在长时间的不断劝告和训导之后，从此以后是他们而不是我们，必须向天主交账。\par

无论如何，我很想知道，你们是否继续劝告你们的弟兄，以及他们是否始终处于同样的懈怠状态；否则我决不会麻烦你们。但事实是，我担心他们会因你们的疏忽和漠不关心而得不到纠正。因为一个人若不断得到劝告和教导，是不可能不变得更好、更勤勉的。我即将引用的谚语固然是常见的，但它证实了这一真理。因为\Quote{水滴不断}，它说，\Quote{能穿石}，水有什么比这更柔软呢？石头又有什么比这更坚硬呢？然而持续的作用能克服自然；若它能克服自然，就更能克服人的意志。亲爱的，基督徒的生活不是儿戏，不是次要的事。我不断说这些事，却毫无效果。\par

2. 你们以为，当我想到在节庆日聚集的群众犹如浩瀚大海，而如今连那群众中最小的一部分也未曾聚集在这里时，我是多么忧伤？那些在节庆日以他们的人数压迫我们的人如今在哪里？我寻找他们，并为他们的缘故而忧伤，因为我察觉有多少走在救恩道路上的人正在丧亡，我蒙受了多么大的弟兄损失，有多少人未能触及救恩之事，而教会的身体大部分犹如一具死寂不动的尸体。你们说：\Quote{这事与我们何干？}关系极大，如果你们不关心你们的弟兄，不劝诫和指导他们，不对他们施加约束，不强行把他们拉到这里，把他们从深沉的懈怠中领出来。因为人不应该只对自己有用，也当对许多人有用，基督清楚地宣告了这一点，当祂称我们为盐\BibleRef{玛 5:13}、酵母和光\BibleRef{玛 5:14}时；因为这些事物对他人有益处和用处。因为灯不是为了自己发光，而是为了那些坐在黑暗中的人；你们是一盏灯，不是为了让你们自己享受光，而是为了领回那边迷路的人。因为如果灯不能照亮坐在黑暗中的人，有何益处？基督徒若对任何人无益，也不领任何人归向德行，又有何益处？同样，盐不是收敛自己，而是坚固那些腐烂的身体部分，防止它们崩解和朽坏。同样，你们也当如此，既然天主指定你们作精神的盐，就要捆扎和坚固那朽坏的肢体，即那些懈怠和污秽的弟兄，将他们从懈怠中解放出来，如同从某种腐败中，使他们与教会身体的其他部分联合。这就是祂称你们为酵母的原因：因为酵母也不发酵自己，而是，虽然它很小，却影响整个面团，无论它有多大。你们也当如此：虽然你们人数很少，但在信德和向天主的热情上，要变得多而有力。因此，正如酵母不因它的微小而软弱，而是因其内在的热和本性的力量而胜出，同样，你们也将能够，如果你们愿意，把比你们自己多得多的数目，领回与你们自己同样的热情程度。现在如果他们以夏季为借口：因为我听到他们说这样的话：\Quote{现在的闷热过分了，灼热的太阳无法忍受，我们不能忍受在人群中被践踏和挤压，浑身冒汗，被炎热和狭小空间压抑。}我为他们感到羞耻，请相信我：因为这些借口是妇人般的；的确，即使对于身体更柔弱、本性更软弱的她们，这样的借口也不足以作为辩护。尽管如此，即使回应这种辩护似乎是一种耻辱，但仍然是必要的。因为如果他们提出这样的借口而毫不脸红，那么我们更不应羞于回应这些事。那么我该对那些提出这些借口的人说什么呢？我要提醒他们关于火窑和火焰中的三位青年，当他们看到火焰从四面环绕他们，包围他们的口、眼睛甚至呼吸时，他们没有停止与万物一同向天主唱那神圣而奥秘的赞美诗，而是站在火堆中间，向万有的共同主献上赞美之歌，其喜乐胜过那些住在百花盛开的田野中的人。并且与这三位青年一起，我认为也应当提醒他们关于巴比伦的狮子，以及达尼尔和狮子坑\BibleRef{达 6:24}。不仅这一个，还有另一个坑，以及耶肋米亚先知，和他被淤泥淹没至颈项的事\BibleRef{耶 38:5}。从这些坑中出来，我要把这些以炎热为借口的人领到监狱，向他们展示那里的保禄和息拉，他们被牢牢地锁在足枷中，浑身青肿和伤痕，遍体鳞伤，却仍在半夜向天主唱赞美诗，举行神圣的守夜。难道这不奇怪吗？那些圣人们，无论是在火窑和火焰中，在坑中，在野兽中，在淤泥中，在监狱和足枷中，在鞭伤和狱卒以及难以忍受的痛苦中，从不抱怨这些事情，反而不断地以极大的精力和热忱念诵祈祷和圣歌，而我们这些从未经历过他们无数苦难中无论大小的人，却因炙热的太阳和一点短暂的炎热辛劳而忽略自己的救恩，离弃聚会而走散，去参加那些完全不健康的集会来败坏自己？当神圣谕言的甘露如此丰沛时，你们却以炎热为借口吗？基督说：\BibleQuote{我所赐给他的水，要在他内成为湧到永生的水泉；}\BibleRef{若 4:14} 又说：\BibleQuote{信从我的人，就如经上所说，从他的腹中要流出活水的江河。}\BibleRef{若 7:38} 请告诉我：当你们有了精神的水井和河流，你们还害怕物质的炎热吗？如今在市场上有那么多的喧嚣拥挤和灼热的风，你们为何不以窒息和炎热为借口而不去呢？因为你们不可能说在那里能享受更凉爽的温度，而所有热气都集中在你们这里；事实正好相反：这里因着铺石地面和建筑的其他构造（因为它被建得很高），空气更轻盈凉爽；而在那里，阳光从各方面都很强烈，且有许多拥挤、蒸汽、尘土以及其他比这些都更增加不适的事物。由此可知，这些无谓的借口是懈怠和懒惰性情的产物，缺乏圣神的火。\par

3. 我这些话并不那么针对他们，而是针对你们——你们没有带他们来，没有唤醒他们脱离懒散，没有吸引他们到这救恩的筵席。家中的奴隶当他们要共同完成某项任务时，会召集同伴；但你们，当你们要聚集举行这属神的礼仪时，却因疏忽而使同仆们被剥夺了益处。你们说：\Quote{但他们若不想来呢？}你们要藉着不断的催迫使祂们渴望：因为如果他们看到你们坚持，他们必定会渴望。不，这些只是借口和托词。这里有多少父亲没有让自己的儿子同他们站在一起？带领你们的一些儿女到这里的难处在哪里？由此可见，所有其他留在外面的人的缺席，不仅由于他们自己的懒散，也由于你们的疏忽。现在至少，如果从前没有，你们要振作起来，让每人带着一位家人进入圣堂：让他们互相激励、催促前来参加这里的集会，父亲带着儿子，儿子带着父亲，丈夫带着妻子，妻子带着丈夫，主人带着奴隶，兄弟带着兄弟，朋友带着朋友；或者更确切地说，我们不仅邀请朋友，也要邀请仇敌到这共同的宝藏来。如果你的仇敌看到你关心他的益处，他无疑会放下仇恨。\par

你对他说：\Quote{难道你不羞愧、不脸红吗？你看犹太人何等严谨地遵守安息日，从傍晚起就停止一切工作。若他们在预备日看到太阳西斜，便中断事务，缩短买卖。若有人在傍晚前向他们买了东西，到傍晚才来付钱，他们决不会收取，也不会接受那钱。}我何必提市场货价和商业交易呢？即使能得着宝藏，他们也宁愿失去收益，而不愿违犯他们的法律。犹太人既然如此严谨——况且他们是在不合时宜地持守法律，固守一种无益反而有害的仪式——而你这超越影相、蒙恩得见正义的太阳、被列为天上国民的人，难道不会彰显同样的热忱吗？那些不合时宜地固守错误的人尚且如此，你这受托付于真理的人，即使只被召来这里短短的一段时间，也不能把它花在听天主的圣言上吗？请问你能得到什么宽恕？你能提出什么合理公正的答复？一个如此冷漠懒散的人，断不能获得宽恕，即使他把世俗事务的需要说上万遍作为借口。难道你不知道，如果你来朝拜天主并参与这里的礼仪，手头的事务会变得容易得多吗？你有世俗的忧虑吗？正因如此，你更该来这里，藉着在此的时光为自己赢得天主的恩宠，然后安心离去；使祂成为你的盟友，使你因天上之手的援助而成为恶魔不能战胜的。如果你受益于司铎们的祈祷，参与共同的祈祷，聆听天主的圣言，为自己赢得天主的助佑，然后全副武装地出发，那么连魔鬼本身也再不能直视你，更何况那些渴望侮辱和毁谤你的恶人呢？但如果你从家直接走到市场，没有这些武器，你就会轻易被所有侮辱你的人制伏。这就是为什么无论在公事还是私事中，许多事情出乎意料地发生：因为我们没有首先勤务属神的事，其次才是世俗的事，而是颠倒了次序。因此，事物应有的顺序和正确的安排被打乱了，我们的一切事务充满了混乱。你能想象我观察到这些时是多么痛苦和忧伤吗？每当公共圣日和庆节临近，全城的居民都聚集而来，虽无人召唤；但圣日和庆节过后，即使我们整天不停地呼喊，也没有人理会。当我反复思量这些事时，常深深叹息，对自己说：劝告或建议有何用呢？你们只是出于习惯行事，并不因我的教导而变得更热心。因为在庆节你们不需要我的劝告，而它们过后，你们从我的教导中毫无所得；你们岂不是证明，就你们而言，我的讲论是多余的吗？\par

4. 也许许多听见这些话的人感到忧伤。但懒散的人没有这种感受：否则他们就会抛弃疏忽，就像我们每天为你们的事焦虑一样。你们从世俗的交易中获得了多少益处，而与此相比，你们遭受了多少损失？从任何其他集会或聚集离开，都不可能像从此地度过的时间那样获得如此多的益处——无论是法庭、议会，甚至是皇宫本身。因为我们并不将治理国家、城市或统帅军队的职责交给进入这里的人，而是另一种比帝国本身更尊贵的治理；或者不如说，我们并不亲自交付，而是由圣神的恩宠交付。\par

那么，比帝国统治更尊贵的统治是什么，那些进入这里的人所领受的又是什么？他们被训练去控制不顺从的情欲，管理邪恶的贪念，驾驭愤怒，调节恶意，制服虚荣。皇帝坐在御座上，头戴王冠，并不像那个人那样尊贵，他将自己内在的正确理性提升到管理卑劣情欲的宝座，并通过对其的统治，如同将一顶荣耀的王冠系在他的额上。试问，当灵魂被情欲所俘虏时，紫色的袍子、金线织成的衣服和镶有宝石的冠冕有什么益处呢？当我们内在的统治元素陷入可耻可怜的奴役状态时，外在的自由又有何益处？因为正如当高烧深入内部，灼烧所有内脏时，即使身体外表没有以同样方式受影响，也无法从中得到任何益处；同样，当我们的灵魂被内在的情欲猛烈挟持时，任何外在的统治，甚至帝国宝座，都没有益处，因为理性被情欲的暴力篡位从帝国的宝座上废黜，并在它们的叛乱运动下弯腰颤抖。为了防止这种情况发生，先知和宗徒们从各方面共同帮助我们，压制我们的情欲，驱逐我们内在非理性元素的一切野蛮，并将一种远比帝国统治更尊贵的统治方式托付给我们。这就是为什么我说，那些剥夺自己这种关怀的人，会在要害处受到打击，承受比任何其他方面所能造成的更大损害，因为那些来到这里的人得到的好处比他们从任何其他来源所能得到的都要大：正如圣经所宣告的。法律说：\BibleQuote{你不可空手出现在主面前；}\BibleRef{出 23:15} 也就是说，不可没有祭献就进入圣殿。如果没有祭献就不该进入天主的殿，那么我们更应该带着我们的弟兄一起进入集会：因为这种祭献和奉献比那一种更好，当你把一个灵魂带进教会时。你看见过那些受过训练的鸽子吗？当它们被放出时，如何寻找其他的鸽子？让我们也这样做。因为如果无理性的动物能够寻找同类的动物，而我们这些被赋予理性和如此多智慧的人却忽略了这种追求，我们将有什么借口呢？我在前一次讲道中用这些话劝告你们：\Quote{你们每个人去邻居的家里，等他们出来，抓住他们，带领他们到他们的共同母亲那里：效法那些对去剧院着迷的人，他们勤奋地安排互相见面，所以在黎明时分就等待去观看那邪恶的表演。}然而我并没有通过这个劝告取得任何效果。因此我再次发言，并且不会停止发言，直到我说服了你们。听到没有益处，除非伴随着实践。它使我们的刑罚更重，如果我们不断听到同样的事情，却不做任何事情。刑罚会更重，请听基督的声明：\BibleQuote{我若没有来对他们说话，他们就没有罪；但如今他们没有借口掩饰他们的罪了。}\BibleRef{若 15:22} 宗徒也说：\BibleQuote{因为听法律的，不将在祂面前成义。}\BibleRef{罗 2:13} 这些话是对听众说的；但当祂要教导讲者，即使他也不会从他的教导中获得任何益处，除非他的行为与他的教义密切相符，他的生活方式与他的言论和谐，请听宗徒和先知是如何对他说的：后者说：\Quote{但天主对罪人说：你为什么传讲我的律法，把我的盟约放在你口中，而你却憎恨管教？}宗徒对那些同样认为自己的教导很了不起的人这样说：\BibleQuote{你自信是瞎子的向导，黑暗中人的光，愚昧人的导师，婴孩的教师：所以你教导别人，却不教导自己吗？}\BibleRef{罗 2:19} 既然它既不能有益于我这讲者说话，也不能有益于你们听者听讲，除非我们遵从所说的话，反而会增加我们的定罪，让我们不要将我们热忱的表现仅限于听讲，而要从行为上遵守所说的话。因为持续聆听神圣之言确实是一件好事：但当听讲所带来的好处没有与之相连时，这件好事就变得无用了。\par

因此，为使你们不致徒然聚集在此，我将不懈地恳切劝勉你们，如我以前多次恳劝你们那样：\Quote{带领你们的弟兄到我们这里来，劝诫那些流浪者，不仅用言语，更要用行为劝告他们。} 这是更有力的教导——即通过我们的举止和品行所传达的——即使你一言不发，但在你离开此集会之后，通过你的神态、目光、声音以及你其余的全部行为，向那些留在后面的人展示你所获得的益处，这就足以劝勉和告诫。\newline{} 因为我们应当从这地方出去，如同从某神圣圣殿出去一样，如同从天上降下的人，变得沉稳、有智慧，凡事举止言谈合宜；当妻子看见丈夫从集会归来，父亲看见儿子，朋友看见朋友，仇人看见仇人，愿他们都感受到你由此处获得的益处；如果他们觉察到你变得更加温和、更有智慧、更虔诚，他们就会感受到。\newline{} 想一想你们这些已受启蒙领受奥秘的人享有何等的特权，与何等团体一起奉献那神秘的赞美诗，与何等团体一起高声喊出\Quote{\OriginalTerm{三圣颂}{Ter sanctus}}。教导\Quote{那些在外的}，说你们已加入色辣芬的合唱，已列于天上邦国的公民之列，已登记在天使的歌团中，已与主交谈，已与基督为伴。\newline{} 如果我们这样规范自己，当我们出去面对那些留在后面的人时，我们就无需说什么了；然而因我们的受益，他们将会感知到自己的缺失，并会急忙前来，以便自己也能享受同样的益处。因为当他们仅凭感官看见你灵魂的美丽闪耀时，即使他们是愚笨至极的人，也会爱上你的美好外貌。因为如果形体之美能激发观看者，那么灵魂的匀称更能感动观看者，并激励他燃起同样的热忱。因此，让我们装饰我们内在的人，当我们出去时，要记念这里所说的话：因为尤其在那里是记念它们的适当时候；正如运动员在赛场上展示他在训练学校中所学的一切一样，我们也应在外面世界的交往中展示我们在这里所听到的一切。\par

5. 因此，要牢记这里所说的话，以便当你出去后，魔鬼藉着愤怒或虚荣，或任何其它情欲抓住你时，你能回忆起在这里所受的教导，并能轻松摆脱那恶者的掌控。\newline{} 难道你们没有看见在训练场上的角力师傅们，他们经过无数竞赛后因年迈而免除角力，坐在场地界线外的尘埃旁，向场内的角力者呼喊，告诉一个人抓对方的手，或拖对方的腿，或抓住后背，并用许多诸如此类的指示，说：\Quote{你若如此这般，就能轻易摔倒对手。} 他们对其弟子有极大的帮助吗？\newline{} 同样，你要仰望你的训练师傅——蒙福的保禄，他在无数胜利之后，现在坐在界限外，我是指此现世生命，他向正在角力的我们大声呼喊，藉着书信向我们高喊，当他看见我们被愤怒和怨恨所胜，被情欲所窒息时：\BibleQuote{如果你的仇人饿了，就给他吃；如果他渴了，就给他喝；} \BibleRef{罗 12:20}——一条美丽的诫命，充满属神的智慧，对行动者和领受者都有益处。\newline{} 但这段话的提醒引起了许多困惑，似乎与先前说话者的情感不一致。这困惑的性质是什么？就是那句话：\BibleQuote{你这样行，就是把炭火堆在他的头上。} 因为藉着这些话，他对行善者和领受者都造成了伤害：对后者，是点着他的头，把炭火堆在其上；因为他从接受食物和饮料中能得到什么益处，相比之下他因头上堆炭火而受的恶呢？如此，领受恩惠者受了冤屈，更重的报复被加在他身上；而行善者也以另一种方式受了伤害。因为当他怀着报复的期望而行善时，他从中能得到什么？因为那为把炭火堆在仇人头上而给他食物和饮料的人，不会变得仁慈和善良，反而残忍和严厉，因小恩施以重罚。还有什么比为了把炭火堆在别人头上而养活他更不仁慈的呢？这就是矛盾：现在需要加上解释，以便藉着那些似乎违反律法字句的事，你能清楚看见立法者的一切智慧。那么，解释是什么呢？\par

那位伟大而高尚的人深知，与仇敌迅速和好是一件沉重而困难的事；沉重而困难，并非因事本身，而是因我们道德上的懈怠。但他不仅命令我们与仇敌和好，还要喂养他；这比前者更为沉重。因为若有人仅因见到激怒自己的人就怒不可遏，又怎会愿意在仇敌饥饿时喂养他呢？我何必谈及见到他们就发怒呢？若有人在人群中提及这些人，仅仅说出他们的名字，就会在我们的想象中重燃创伤，加剧情感的狂热。\Person{保禄}深知这一切，希望使那难以纠正的事变得平滑容易，并说服那不能忍受见到仇敌的人，预备好施行上述的恩惠，于是加上了关于火炭的话，为使人因盼望报复而赶紧向那得罪自己的人施行此善举。正如渔夫将钩子四面用饵包裹，呈给鱼群，以便其中一条急于取食而被钩住，轻易被擒：同样，\Person{保禄}也愿意引导那受亏负的人向亏负他的人施恩，他不向他呈现属灵智慧的赤裸钩子，而是好像用某种饵包裹起来，我是指那\Quote{火炭}，邀请那受侮辱的人，因指望惩罚，而向那得罪自己的人施此善举；但当他前来，便从此将他牢牢掌握，不容他逃脱，因为行为本身的性质将他与仇敌联结；他几乎对他说：\Quote{你若不愿因虔敬而喂养那得罪你的人，至少因盼望惩罚而喂养他。}因为他知道，人一旦着手施行此善举，便有了开端，和好的道路便为他敞开。因为确实，没有人会忍心一直视他给自己饮食的人为仇敌，即使最初他是出于报复的盼望而行此事。因为时间流逝，会缓解他愤怒的紧张。正如渔夫若呈现赤裸的钩，绝不会引诱鱼，但当他用饵包裹，鱼在不觉中将其吞入口中：同样，\Person{保禄}若没有提出惩罚的指望，就绝不会说服那些受亏负的人去善待那些得罪他们的人。因此，他想要说服那些厌恶退缩、甚至见到仇敌就瘫痪的人，向他们施行最大的恩惠，便提及了火炭，并非要将这些人投入无情的惩罚，而是为使他先说服受亏负的人因盼望惩罚而善待仇敌，以后能逐渐说服他们完全放弃愤怒。\par

6. 这样，他鼓励了那受亏负的人；但也要观察他如何将那作恶的人再次与那被激怒的人联合。首先是通过恩惠的方式本身：（因为没有人如此卑劣无情，以至于领受饮食后，仍不愿成为供养他之人的仆人和朋友）；其次是通过对报复的恐惧。因为这段经文：\BibleQuote{你这样行，就是把火炭堆在他头上}似乎确实是对那施食者说的；但它更特别触及那惹事的人，为使他因惧怕这惩罚而不致持续处于敌对状态，并知道若他始终怀有敌意，领受饮食可能给他带来极大的祸害，从而抑制他的愤怒。因为这样他就能熄灭那火炭。因此，所提出的惩罚和报复既促使受亏负的人去善待那得罪他的人，又阻止和约束那挑衅者，并推动他与那给他饮食的人和好。\Person{保禄}因此用双重的纽带将两人联结：一条倚靠恩惠，另一条倚靠报应行为。因为困难在于为和好开辟一个开端和入口；但是一旦以任何方式清理了这一点，此后的一切就平坦容易了。因为即使起初那被惹怒的人因盼望惩罚而喂养他的仇敌，但通过给他食物这一行为成为他的朋友后，他就能驱除报复的欲望。因为一旦成为朋友，他就不再有这种期望而去喂养那已与他和好的人。同样，那挑衅者，当他看见受亏负的人选择给他饮食时，他会抛开一切敌意，既因为这一行为，也因为他惧怕那为他存留的惩罚，即使他极其强硬、粗暴、铁石心肠，也会因施食者的仁慈而羞愧，并因惧怕接受食物后若继续为敌所保留的惩罚而感动。\par

为此，\Person{保禄}在劝勉中并未止步于此，而是在使双方都清除了义怒之后，进而矫正他们的心态，说道：\Quote{不可为恶所胜。}因为，\Quote{若你们继续怀怨并寻求报复，似乎是在战胜你的仇敌，但实际上你是被恶所胜，即被义怒所胜；所以若你愿意得胜，就该和好，不要攻击你的对手；}那才是辉煌的胜利，在其中藉着善，即藉着容忍，你克服了恶，驱散了义怒和怨恨。但受伤害的人，若被激情点燃，是不会忍受这些话的。因此，当\Person{保禄}满足了他的义怒后，便引导他走向和好的最佳理由，不让他长久停留在报复的邪恶希望中。你看到这位立法者的智慧了吗？为使你知道他引入这条法律只是为那些否则不愿彼此和解之人的软弱起见，请听\Person{基督}在就同一主题制定法律时，是如何提出与宗徒不同的赏报的：祂说\Quote{爱你们的仇敌，善待恨你们的人}，即给他们食物和饮料，却没有加上\Quote{因为你这样做，就是将炭火堆在他头上}；而是说什么？\Quote{好使你们成为你们在天之父的子女。}\BibleRef{玛 5:44}自然而然地，因为祂是在向\Person{伯多禄}、\Person{雅各伯}、\Person{若望}以及其余宗徒团体讲论，所以祂提出了那样的赏报。但若你说即使这样理解，这条诫命仍是沉重的，那么你既又一次为我为\Person{保禄}所作的辩护加分，却也剥夺了你一切宽容的借口。因为我能向你证明，这在你看来沉重的事，在旧约时代已经实现了，那时属灵智慧的光照不如现在这般大。为此，\Person{保禄}也不是用自己的话引入这条法律，而是藉用了最初引入它的人所用的确切词语，为使那些不遵守它的人无可推诿：因为\Quote{若你的仇敌饿了，就给他吃；若他渴了，就给他喝}\BibleRef{箴 25:21}这条诫命，最初不是\Person{保禄}的话，而是\Person{撒罗满}的话。为此他引用了这些话，为使听者相信，一个被提升到如此高度智慧的人，却将一条古人常常遵行的旧法律视为沉重和难堪，这是最可耻的事之一。你问，古人中谁遵行了它？有许多人，但尤其是\Person{达味}，更为丰盛地做到了。他不仅给予他的仇敌食物或饮料，还多次在他遇险时救他脱离死亡；当他有机会杀他时，他放过了他一次、两次，甚至多次。至于\Person{撒乌耳}，在\Person{达味}为他做了数不清的善事、辉煌的胜利以及关于\Person{哥肋雅}所行的拯救之后，他如此仇恨和厌恶他，甚至不忍心用他的名字称呼他，而是以他父亲的名字称呼他。因为有一次，节期临近，\Person{撒乌耳}对他设下了奸计和残酷的阴谋，却没有见他到来——\Quote{\Person{叶瑟}的儿子在哪里？}\BibleRef{撒上 20:23}他叫他以他父亲的名字，既因为出于仇恨他无法忍受对他本名的回忆，也因为想借提及他低微的出身来损害那位义人的崇高地位——这是一种可悲且可鄙的想法：因为确实，即使他对父亲有一些指控，也丝毫不能伤害\Person{达味}。因为每人要为自己的行为负责，并以此受赞美或指责。但事实上，因无可指责的恶行可说，他便提出他低微的出身，指望借此遮蔽他的荣耀，这其实是愚蠢之极。因为出身卑微谦逊之人，即\Quote{\Person{叶瑟}的儿子}，算是什么过错呢？但当\Person{达味}发现他睡在洞内时，并没有称他为\Quote{\Person{克士}的儿子}，而是以他的尊称：\Quote{因为我决不愿伸手加害上主的受傅者。}\BibleRef{撒上 26:11}他是如此彻底地脱离了义怒和怨恨：他称那对他行了如此大恶、渴求他流血、在他无数善行之后多次企图毁灭他的人为\Quote{上主的受傅者}。因为他不是看\Person{撒乌耳}应受何种对待，而是看他自己应当做什么和说什么，这是道德智慧的最高境界。何以见得？当你的仇敌被关在监狱里，被双重，甚至三重锁链——空间的限制、缺乏援助、睡眠的必要——束缚得牢牢时，难道你还不要求他受到惩罚和报应吗？\Quote{不，}他说：\Quote{因为我现在不是看他该受什么痛苦，而是看我应当做什么。}他不看杀戮的便利，而看对自己合宜的道德智慧的精确遵守。然而，当时的环境中有哪一样不足以促使他动手杀人呢？难道他的仇敌被捆绑交到他手中，不是一个充分的诱因吗？因为我想你知道，当事情方便时，我们更急切地行事，成功的希望增加了我们行动的欲望，这正是当时在他身上发生的事。\par

好啊！那位当时劝告并催促他行事的将领，难道过去的记忆促使他杀戮了吗？这一切没有一样打动他：实际上，正是轻易杀戮的可能性使他避免了这种行为，因为他意识到，天主把撒乌耳交在他手中，正是为了给他提供充分的理由和机会来操练道德智慧。那么，你或许佩服他，因为他没有记住自己过去所受的任何伤害；但我更因另一个原因而对他感到惊讶。那是什么呢？那就是对未来的恐惧没有驱使他向敌人动武。因为他清楚地知道，如果撒乌耳从他手中逃脱，他将会再次成为他的敌人；但他宁愿释放那个曾冤枉他的人而将自己置于危险之中，也不愿通过向敌人动武来保障自己的安全。那么，还有什么能与这个人的伟大和慷慨精神相比呢？当法律命令以眼还眼，以牙还牙，以及同等的报复时，\BibleRef{申 19:21}他不仅没有这样做，反而展现了更大程度的道德智慧？至少，如果他在那时杀了撒乌耳，他仍会保持道德智慧的声誉不受损害，不仅因为他是在自卫，不是首先施暴的人，而且因为他极大的节制使他超越了\Quote{以眼还眼}的诫命。因为他不会以一次杀戮回报一次杀戮；相反，撒乌耳多次试图杀害他（不是一次两次，而是许多次），而达味只会让撒乌耳死一次；不仅如此，如果他出于对未来的恐惧而报复，加上已提到的种种，这也会使他获得毫无折扣的忍耐之赏报。因为那个为自己所受的伤害而发怒并寻求满足的人，无法得到忍耐的赞美；但一个人若抛开过去所有伤害的考虑（尽管它们又多又痛苦），而被迫出于对未来的恐惧并为了保障自己的安全而采取自卫措施，没有人会剥夺他节制的赏报。\par

7. 然而，达味甚至没有这样做，而是发现了一种新奇而独特的道德智慧形式：过去的记忆、对未来的恐惧、将领的煽动、地方的荒僻、杀戮的轻易，以及其他一切，都没有促使他杀人；相反，他饶恕了那个既是他的敌人又曾使他痛苦的人，就好像那人是他的恩人、曾为他做了许多好事一样。那么，我们将会得到怎样的宽恕呢？我们若铭记过去的过犯，报复那些使我们痛苦的人，而那位无辜的人，曾忍受如此巨大的苦难，并预期因饶恕他的敌人而会遭遇更多更糟的祸患，却仍饶恕了他，以至于宁愿自己陷入危险、生活在恐惧战栗之中，也不愿将那个会给他带来无穷麻烦的人处以公正的死刑？\par

我们从此可以识得他的道德智慧，不仅在于他没有杀死撒乌耳——当时他有如此强烈的冲动——也在于他没有对撒乌耳说过一句不敬的话，尽管那位受辱者不会听见他的话。我们常背后说朋友的坏话，而他反而不说那如此亏待他的敌人的坏话。他的道德智慧，我们由此可知；但他的慈爱与温柔关怀，则从他以后所行的事上可以看出。因为他割下撒乌耳衣袍的衣边，又拿走水囊后，远远退去，站着呼喊，并向那位他保全性命的人展示这些，不是出于炫耀和显摆，而是希望藉行为使他相信，自己无故把他当作敌人，并意在赢得他的友谊。然而，即或这样仍不能说服他，而且自己本可下手抓他，他却再次选择离乡背井、客居异地，每日为谋取必要食物而受苦，也不愿留在家里与对手为难。有谁的胸怀比这更仁慈？他实在有理由说：\Quote{主，记念达味和他的一切温良。}我们也当效法他，对仇敌不说也不做恶事，反要量力帮助他们：因为我们对他们的好处，更大过对自己的好处。因为我们被告知：\Quote{因为如果你们宽恕你们的仇敌，}我们被告知\Quote{你们必蒙宽恕。}（\BibleRef{玛 6:14}）宽恕卑劣的过犯，好使你们为自己的过犯获得君王的赦免；但若有人大大亏负了你们，你们所宽恕的亏负越大，你们将获得的赦免也越大。因此，我们受教导要说：\Quote{宽恕我们，如同我们宽恕别人一样}，好使我们明白，我们宽恕的尺度首先是从我们自己开始的。所以，仇敌加给我们的恶越大，他所赐予我们的益处也越大。我们要认真又迫切地与激怒我们的人和好，无论他们的怒气有理无理。因为你们若在此地和好，就能免于来世的审判；但若在仇恨持续期间，死亡突然插入并带走仇恨，那案件的审理必然要带到来世。正如许多人彼此争执时，若在法庭外达成友好谅解，便免去了损失、惊吓和许多风险——案件的结局取决于各人的心意；但若他们各自将事情交托法官，其结果只会是金钱的损失，在许多情况下还要加上惩罚，且仇恨永远持续。同样，我们若在今生和解，就能免除一切惩罚；但若仍存仇恨而离开，进入来世那可畏的审判台前，就必在审判官的判决下受最重的刑罚，且两人都要经受无情的惩罚：那无故发怒的，因为他无理；而那有理发怒的，也因为他虽有理却仍怀恨。因为即使我们受了无理的亏待，也应当宽恕那些得罪我们的人。请注意，祂如何敦促和激励那些无故使人痛苦的人，与那被他们得罪之人和好。\Quote{所以，若你在祭台前献礼物时，在那里想起你的弟兄有什么怨你的事，就把礼物留在祭台前，先去与你的弟兄和好。}（\BibleRef{玛 5:23}）祂没有说：\Quote{集合起来献祭}，而是说\Quote{和好，然后再奉献}。祂说：把礼物留在那里，好让献祭的必要迫使那有理发怒的人也违背本意来和解。请看，当祂说\Quote{宽免欠你们债的人，好使你们的父也宽免你们的过犯}时，祂又如何催促我们去寻找那激怒我们的人。因为祂提出的赏报并不微小，而是远超过成就的伟业。既然如此，我们想想这一切，数算在此事上所得的报偿，并记住消除罪过并不需要很多劳苦和热忱，就让我们宽恕那些亏负我们的人吧。因为别人藉着禁食、哀恸、祈祷、苦衣和灰土几乎不能成就的事——我指的是洗除自己的罪过——我们却可以不用苦衣、灰土和禁食而轻易做到，只要我们发自内心地除去愤怒，并真诚地宽恕那些得罪我们的人。愿和平与爱情的天主，从我们灵魂中驱除一切忿怒、苦毒和怒气，恩准我们彼此紧密相连，按照各部分的适当配合（\BibleRef{弗 4:16}），一心一口，一心一魂，不断向祂献上我们应得的感恩颂歌：因为光荣与权能归于祂，直到永远。阿们。\par

\SourceRecord{Translated by W.R.W. Stephens.；Nicene and Post-Nicene Fathers, First Series；Vol. 9.；Edited by Philip Schaff.；Buffalo, NY: Christian Literature Publishing Co.,；1889.；<http://www.newadvent.org/fathers/1912.htm>.}
